\section{Panel de la sesión matutina: modelo, harness, recompensa y entorno, ¿qué evoluciona primero?}

El panel de la sesión matutina reunió a Yang Mengyue (杨梦月), Zhang Shaokun (张少坤), Qian Cheng (钱成), Dong Guanting (董冠霆) y Hu Mengkang (胡梦康) para discutir el conflicto más directo entre las cuatro presentaciones anteriores: los efectos a corto plazo suelen provenir de la optimización del harness y del context, mientras que la capacidad a largo plazo requiere una actualización del modelo; el Agent necesita diseñar sus propias tareas y su propia recompensa, pero la autoevaluación es también lo que más fácilmente produce reward hacking; las aplicaciones reales exigen exploración abierta, pero a la vez deben tener límites verificables.

\subsection{Por qué la auto-investigación se convirtió en el terreno de prueba de la automejora}

Los panelistas coincidieron en general en que la investigación automatizada es un escenario de alto valor, porque las tareas de investigación incluyen descubrimiento de problemas, búsqueda, experimentación, ejecución de código, interpretación de resultados y generación de reportes, lo que forma de manera natural una retroalimentación de ciclo largo. Un paso más allá, el «auto-research for auto-research» hace que el Agent no solo complete experimentos, sino que también intente mejorar el propio flujo de investigación, el benchmark o sus propias herramientas.

\begin{importantbox}{Consenso del panel: la automejora debe cambiar la conducta futura}
Sin importar si lo que se actualiza es el prompt, el harness, la memoria o los parámetros del modelo, debe observarse un cambio reproducible en tareas posteriores. Generar solo reflexiones, resúmenes o ideas nuevas, sin cambiar la estrategia de ejecución, no constituye una automejora válida.
\end{importantbox}

\subsection{Evolución del modelo o evolución del harness}

La respuesta de ingeniería que dio el panel no es una disyuntiva, sino un tratamiento por escalas de tiempo. Modificar el harness tiene un costo bajo y una retroalimentación rápida, lo que lo hace adecuado para validar nuevas herramientas, nuevos flujos y nuevas restricciones; una vez que cierta estructura se muestra establemente eficaz en una gran cantidad de tareas, se considera interiorizarla como capacidad del modelo mediante entrenamiento. Esto permite iterar rápido sin tener que reentrenar el modelo base cada vez que cambia el producto.

\begin{knowledgebox}{La ruta de la externalización a la interiorización}
Una ruta habitual es: primero implementar la capacidad con un prompt o un workflow, recolectar trayectorias reales y validar el beneficio; después destilar la conducta estable o incorporarla al modelo mediante RL; por último, simplificar el harness. El harness es así, a la vez, un componente del producto y un generador de datos para explorar futuras capacidades del modelo.
\end{knowledgebox}

Los docentes también advirtieron que el harness y el modelo no son completamente ortogonales. El harness cambia la distribución de observaciones, acciones y retroalimentación, y el modelo entrenado, a su vez, cambia cuál es el harness óptimo. Un sistema verdaderamente maduro necesita versionarse en conjunto y evaluarse de manera conjunta.

\subsection{¿Puede el Agent diseñar su propia recompensa?}

Las tareas abiertas rara vez tienen una recompensa continua y sin ambigüedad. Si se deja que el Agent genere su propio proxy reward, este tenderá a optimizar el indicador más fácil de medir, en lugar del objetivo que realmente le importa a la persona. El panel consideró esto el riesgo central de la self-evolution: si la generación de tareas, la respuesta y la calificación quedan bajo el control del mismo modelo, el sistema puede confirmar continuamente su propio progreso dentro de un ciclo cerrado.

\begin{warningbox}{Tres orígenes del reward hacking}
Primero, el indicador proxy omite el objetivo real; segundo, el verificador presenta vulnerabilidades explotables; tercero, el evaluador y la política derivan juntos. Se debe procurar introducir ejecución de programas, estado del entorno, muestreo humano y validación cruzada con múltiples evaluadores, además de conservar pruebas de regresión sobre tareas de versiones anteriores.
\end{warningbox}

\subsection{Entorno, modelo del mundo y exploración activa}

El modelo causal del mundo enfatiza que el Agent debe buscar activamente intervenciones que permitan distinguir entre hipótesis, mientras que Agent-World enfatiza la generación continua de nuevos entornos y tareas. El panel conectó ambas ideas: el generador de entornos aporta el espacio de exploración, el world model predice el valor de la información, la rollout infrastructure ejecuta la interacción y el reward/verifier determina si se obtuvo conocimiento nuevo.

\begin{importantbox}{La tensión entre apertura y verificabilidad}
Cuanto más abierta es una tarea, más probable es que produzca nuevas capacidades, pero también es más difícil calificarla de manera automática; cuanto más fácil de verificar es una tarea, más adecuada resulta para RL, pero puede volverse demasiado estrecha. Un currículo práctico debe combinar tareas con verificación fuerte, tareas de exploración con verificación débil y muestras de revisión humana, en lugar de elegir solo uno de los extremos.
\end{importantbox}

\subsection{Forma de producto en el corto plazo: una federación de Agents verticales}

Los panelistas consideraron que la forma más realista en el corto plazo no es un único Agent omnipotente, sino múltiples Agents verticales que se llaman entre sí mediante MCP, skills o interfaces de servicio. Dominios como código, análisis de datos, búsqueda e investigación formarán primero unidades de alta confiabilidad, que después se combinarán gradualmente en sistemas más complejos. Esta estructura federada también exige declaraciones explícitas de capacidad, permisos y límites de propagación de fallos.

\subsection{Resumen de la sección}

El panel de la sesión matutina propuso una ruta pragmática: experimentar rápido a nivel de harness, recolectar on-policy experience en entornos verificables, interiorizar los patrones estables en el modelo, y al mismo tiempo hacer crecer el entorno y el currículo junto con la capacidad de la política. Todo mecanismo de automejora debe evitar los ciclos cerrados de autoevaluación, conservar la posibilidad de retroceder a versiones anteriores, y emparejar la exploración abierta con verificación externa.
